% ============================================================
% chapters/01-intro.tex -- פרק 1: מבוא ורקע
% Content only. All design comes from template.cls.
% Commands used here that the template must provide:
%   \en{...}            English text inside Hebrew
%   \LR{...}            left-to-right run (year ranges)
%   keybox environment  highlighted box with a title
% ============================================================

\chapter{מבוא ורקע}
\label{ch:intro}

\section{אוגוסט 410: היום שבו רומא נבזזה}
\label{sec:intro-410}

ב־24 באוגוסט 410 נכנסו הוויזיגותים, בהנהגת המלך אלאריך, אל העיר רומא ובזזו אותה במשך שלושה ימים. היה זה הכיבוש הראשון של העיר בידי אויב זר מזה כשמונה מאות שנה, מאז פשיטת הגאלים במאה הרביעית לפני הספירה. מבחינה צבאית ומדינית, רומא כבר לא הייתה אז בירת הקיסרות: חצר הקיסר במערב ישבה ברוונה, והשלטון האמיתי התפצל מזמן בין מזרח למערב. אבל מבחינה סמלית, הזעזוע היה עצום. העיר שכבשה את העולם נכבשה בעצמה, והידיעה עברה מקצה הים התיכון ועד קצהו.

לאירוע הזה הייתה תפנית אירונית, שהיא גם נקודת המוצא של הדוח הזה: הוויזיגותים עצמם היו נוצרים. הם השתייכו לזרם האריאני, שנחשב לכפירה בעיני הכנסייה הקתולית, אבל הם כיבדו את הכנסיות של רומא כמקומות מקלט ונמנעו מלפגוע בחלק מהמסתתרים בהן \cite{augustine_citygod}. עיר נוצרית נבזזה בידי צבא נוצרי, שלושים שנה בלבד אחרי שהנצרות הוכרזה כדת הרשמית של הקיסרות.

\begin{keybox}{התקופה בקצרה}
\begin{itemize}
  \item \textbf{312} – קונסטנטינוס מנצח בקרב הגשר המילווי ומתחיל לתמוך בנצרות.
  \item \textbf{380} – תאודוסיוס הראשון קובע את הנצרות הניקאית כדת הרשמית של הקיסרות.
  \item \textbf{395} – הקיסרות מתחלקת סופית בין שני בניו של תאודוסיוס: מזרח ומערב.
  \item \textbf{410} – הוויזיגותים בוזזים את רומא.
  \item \textbf{476} – הקיסר האחרון במערב מודח; הקיסרות המזרחית ממשיכה להתקיים.
\end{itemize}
\end{keybox}

\section{ויכוח בן 1,600 שנה}
\label{sec:intro-debate}

השאלה ״האם הנצרות אשמה?״ עלתה כמעט מיד. בני העילית הפגאנית ברומא טענו שהעיר נענשה משום שזנחה את האלים שהגנו עליה במשך מאות שנים. בתגובה כתב אוגוסטינוס, הבישוף של היפו שבצפון אפריקה, את חיבורו הגדול \en{De Civitate Dei} (״עיר האלוהים״), בין השנים \LR{413--426}. אוגוסטינוס טען שגורלן של ממלכות ארציות אינו מבחן לאמת הדתית, ושמקומו האמיתי של הנוצרי הוא ״עיר האלוהים״ ולא עיר האדם \cite{augustine_citygod}.

כאלף ושלוש מאות שנה אחר כך החזיר אדוארד גיבון את הטענה הפגאנית לקדמת הבמה. בספרו \en{The History of the Decline and Fall of the Roman Empire}, שראה אור בשישה כרכים בין \LR{1776--1789}, תיאר גיבון את נפילת רומא כ״ניצחון הברבריות והדת״. לטענתו, הנצרות הסיטה את מבטם של הרומאים מן המדינה אל העולם הבא, משכה אנשים מוכשרים אל הכנסייה במקום אל הצבא והממשל, ופילגה את החברה בוויכוחים תאולוגיים \cite{gibbon_decline}.

ההיסטוריונים של העשורים האחרונים דוחים ברובם את הנצרות כהסבר מרכזי. פיטר הת'ר רואה בלחץ החיצוני של עמי הגבול, שהתחזקו דווקא בזכות מגעם הארוך עם רומא, את הגורם המכריע \cite{heather2005fall}. בריאן וורד־פרקינס מדגיש שהנפילה הייתה קטסטרופה כלכלית וחומרית אמיתית ולא ״שינוי״ שקט \cite{wardperkins2005fall}. קייל הרפר מוסיף את האקלים ואת המגפות כגורמים מרכזיים לצד תהליכים צבאיים, מדיניים וכלכליים \cite{harper2017fate}. ופיטר בראון, מגדולי חוקרי התקופה, מערער ישירות על הטענה שהעושר שזרם אל הכנסייה החליש את יכולתה של רומא לעמוד בפני הפולשים \cite{brown2012needle}.

\begin{keybox}{העובדה שמסבכת את הוויכוח}
גם הקיסרות הרומית במזרח הייתה נוצרית, ובאותה מידה כמו המערב, אבל היא שרדה עוד כאלף שנה, עד \LR{1453}. ובמערב עצמו, הממלכה הפרנקית, שהייתה נוצרית כמעט מיום הקמתה, הפכה לכוח המרכזי באירופה. אם הנצרות החלישה מדינות, איך מסבירים את שני המקרים האלה?
\end{keybox}

\section{למה עכשיו: כלים חדשים לשאלה ישנה}
\label{sec:intro-why-now}

עד לאחרונה, הוויכוח התנהל כמעט כולו בטיעונים איכותניים. מחקרים רבים התמקדו במקרה אחד ובדוגמאות שתמכו בעמדתם, והשוואה שיטתית בין מדינות שונות, באותם מדדים, הייתה קשה מאוד לביצוע. שלוש התפתחויות של השנים האחרונות משנות את המצב:

\begin{enumerate}
  \item \textbf{מאגרי נתונים היסטוריים.} מאגר \en{Seshat: Global History Databank} מקודד באופן שיטתי נתונים על מבנה חברות לאורך ההיסטוריה. מחקר שהתבסס עליו קידד נתונים על 414 חברות מ־30 אזורים ברחבי העולם, לאורך עשרת אלפים השנים האחרונות \cite{turchin2018quantitative}, ומחקר המשך בדק כמותית 17 משתנים שהוצעו בתאוריות שונות כמניעים של מורכבות חברתית \cite{turchin2022disentangling}.
  \item \textbf{ספריות טקסטים דיגיטליות.} רוב המקורות הלטיניים והיווניים של התקופה זמינים היום בפורמט דיגיטלי, וניתנים לחיפוש ולניתוח ממוחשב.
  \item \textbf{מודלי שפה.} מודלים שאומנו על טקסטים לטיניים מסוגלים לנתח שפה עתיקה; מודל \en{Latin BERT}, לדוגמה, אומן על כ־642.7 מיליון מילים לטיניות מהתקופה הקלאסית ועד ימינו \cite{bamman2020latinbert}. מודלי שפה גדולים מאפשרים לסווג אלפי קטעים לפי נושא, בהיקף שאף חוקר יחיד אינו יכול לקרוא.
\end{enumerate}

השילוב בין שלושת הכלים האלה מאפשר לבחון את הוויכוח בן 1,600 השנים גם כשאלה אמפירית: למדוד, להשוות ולבחון השערות, בהיקף ובשיטתיות שהיו קשים להשגה בעבר.

\section{המוטיבציה}
\label{sec:intro-motivation}

לשאלה יש ערך שחורג מעבר להיסטוריה של העת העתיקה. היחסים בין דת למוסדות המדינה עומדים במרכז הוויכוח הציבורי בחברות רבות גם היום, ולעתים קרובות שני הצדדים משתמשים ב״לקחי ההיסטוריה״ – ובמיוחד בנפילת רומא – כדי לתמוך בעמדתם. מחקר שמבחין בין \textbf{הדת עצמה} לבין \textbf{האופן שבו היא משתלבת במוסדות} יכול לתרום לדיון הזה בסיס עובדתי במקום אנלוגיות.

המחקר המוצע בדוח משווה שלוש מדינות נוצריות שהגיעו לתוצאות שונות לגמרי: הקיסרות הרומית המערבית, הקיסרות הרומית המזרחית והממלכה הפרנקית. ההשערה המרכזית היא שהדת כשלעצמה לא הכריעה את גורלן; מה שהכריע הוא מידת ההשתלבות בין הכנסייה לבין מוסדות המדינה.

\section{מבנה הדוח}
\label{sec:intro-structure}

פרק~\ref{ch:problem} מגדיר את הבעיה, מנסח את שאלת המחקר ואת ההשערות, ומציג את מפרט הדרישות של המחקר. פרק~\ref{ch:literature} סוקר את הספרות המחקרית לפי נושאים, מציג את מאמר הדגל ומגדיר את פער הידע. פרק~\ref{ch:market} סוקר את הכלים ומאגרי הנתונים הקיימים לחוקרי התקופה. פרק~\ref{ch:alternatives} משווה שלוש חלופות שיטה ובוחר ביניהן במטריצת החלטה משוקללת. פרק~\ref{ch:workplan} מציג את תוכנית העבודה, ופרק~\ref{ch:risks} את הסיכונים ואת דרכי ההתמודדות איתם. פרק~\ref{ch:summary} מסכם ומציג את הצעדים המיידיים לאחר אישור הדוח.
